\documentclass[a4,11pt]{aleph-notas}

% -- Paquetes adicionales
\usepackage{enumitem}
\usepackage{aleph-comandos}

% -- Datos
\institucion{Facultad de Ciencias Exactas, Naturales y Ambientales}
\carrera{Catálogo STEM}
\asignatura{Matemática III}
\tema{Resumen 06: Derivadas de orden superior}
\autor[A. Merino]{Andrés Merino}
\fecha{Periodo 2026-1}

\logouno[0.16\textwidth]{Logos/logoPUCE_04_ac}
\definecolor{colortext}{HTML}{1957d1}
\definecolor{colordef}{HTML}{1957d1}
\fuente{montserrat}


\begin{document}

\encabezado

En este resumen consideraremos $\Omega\subset \R^n$ un conjunto abierto y $n>1$.

%%%%%%%%%%%%%%%%%%%%%%%%%%%%%%%%%%%%%%
\section{Derivadas de orden superior}
%%%%%%%%%%%%%%%%%%%%%%%%%%%%%%%%%%%%%%

\begin{defi}[Derivadas parciales de orden superior]
    Dados $\func{f}{\Omega}{\R}$ un campo escalar tal que, para $i=1,\ldots, n$, existe $D_if(x)$ para todo $x\in \Omega$, para $j=1,\ldots, n$ y $a\in \Omega$, se define la derivada parcial mixta por
    \[
        D_{j,i} f(a)=\frac{\partial^2 f}{\partial x_j\partial x_i}(a)
        =\frac{\partial f}{\partial x_j}\l( \frac{\partial f}{\partial x_i}\r)(a)
        =D_j(D_i f)(a).
    \]
\end{defi}

\begin{teo}
    Dados $\func{f}{\Omega}{\R}$ un campo escalar tal que sus derivadas parciales mixtas existan y sean continuas en $\Omega$, entonces
    \[
        D_j(D_i f)(a) = D_i(D_j f)(a)
    \]
    para todo $i,j\in\{1,\ldots, n\}$ y $a\in \Omega$.
\end{teo}

\begin{defi}[Matriz hessiana]
    Dado $\func{f}{\Omega}{\R}$ un campo escalar diferenciable con continuidad en $\Omega$, además consideramos $\func{\nabla f}{\Omega}{\R^n}$ como un campo vectorial. Si $\nabla f$ posee matriz jacobiana en $a\in \Omega$, entonces se define la matriz hessiana de $f$ por
    \[
        H_f(a)=
        \begin{bmatrix}
         D_{1,1}f(a) & D_{1,2}f(a) & \cdots & D_{1,n}f(a) \\
         D_{2,1}f(a) & D_{2,2}f(a) & \cdots & D_{2,n}f(a) \\
         \vdots & \vdots &  \ddots & \vdots\\
         D_{n,1}f(a) & D_{n,2}f(a) & \cdots & D_{n,n}f(a)
        \end{bmatrix}.
    \]
\end{defi}

\begin{advertencia}
    La matriz hessiana puede ser entendida, de cierto modo, como el representante de la segunda derivada de la función, la cual también es una aplicación lineal.
\end{advertencia}

\begin{advertencia}
    Si las derivadas parciales mixtas de una función son continuas, entonces su matriz hessiana es simétrica. En este caso, se dice que la función es dos veces diferenciable con continuidad.
\end{advertencia}

\newpage
%%%%%%%%%%%%%%%%%%%%%%%%%%%%%%%%%%%%%%
\section{Fórmulas de Taylor}
%%%%%%%%%%%%%%%%%%%%%%%%%%%%%%%%%%%%%%

\begin{teo}[Fórmula de Taylor de primer orden]
    Sea $\func{f}{\Omega}{\R}$ un campo escalar diferenciable y $a\in \Omega$, entonces existe $r>0$ tal que para todo $h\in\R^n$ que cumple $\|h\|<r$, se tiene que
    \[
        f(a+h) = f(a) + \nabla f(a) \cdot h + \|h\| E(a,h),
    \]
    donde $E(a,h)\to 0$ cuando $h\to 0$.
\end{teo}

\begin{prop}[Aproximación lineal]
    Sea $\func{f}{\Omega}{\R}$ un campo escalar diferenciable y $a\in \Omega$, entonces la función definida por
    \[
        F(x) = f(a) + \nabla f(a) \cdot (x-a),
    \]
    es una buena aproximación lineal de la función $f$ cerca de $a$, es decir,
    \[
        \lim_{x\to a} \frac{F(x)-f(x)}{\|x-a\|} = 0.
    \]
\end{prop}

\begin{teo}[Fórmula de Taylor de segundo orden]
    Sea $\func{f}{\Omega}{\R}$ un campo escalar dos veces diferenciable y $a\in \Omega$, entonces existe $r>0$ tal que para todo $h\in\R^n$ que cumple $\|h\|<r$, se tiene que
    \[
        f(a+h) = f(a) + \nabla f(a) \cdot h +  \frac{1}{2!}[h]^T H_f(a) [h] +\|h\|^2 E(a,h),
    \]
    donde $E(a,h)\to 0$ cuando $h\to 0$.
\end{teo}

\begin{prop}[Aproximación cuadrática]
    Sea $\func{f}{\Omega}{\R}$ un campo escalar dos veces diferenciable y $a\in \Omega$, entonces la función definida por
    \[
        F(x) = f(a) + \nabla f(a) \cdot (x-a) + \frac{1}{2!}[x-a]^T H_f(a) [x-a],
    \]
    es una buena aproximación cuadrática de la función $f$ cerca de $a$, es decir,
    \[
        \lim_{x\to a} \frac{F(x)-f(x)}{\|x-a\|^2} = 0.
    \]
\end{prop}


\end{document}
